% tables/algorithm_comparison.tex — Master algorithm comparison table
% Included via % tables/algorithm_comparison.tex — Master algorithm comparison table
% Included via % tables/algorithm_comparison.tex — Master algorithm comparison table
% Included via % tables/algorithm_comparison.tex — Master algorithm comparison table
% Included via \input{tables/algorithm_comparison} from a section file.
% Designed for landscape orientation; wrap the \input{} call inside
% a \begin{landscape}...\end{landscape} environment (from pdflscape) or
% place it on a standalone landscape page.

\begin{footnotesize}
\setlength{\LTcapwidth}{\linewidth}
\setlength{\tabcolsep}{3.5pt}
\renewcommand{\arraystretch}{1.25}

\begin{longtable}{%
  >{\raggedright\arraybackslash}p{2.1cm}   % Algorithm
  >{\raggedright\arraybackslash}p{1.5cm}   % Family
  >{\raggedright\arraybackslash}p{1.8cm}   % Input Shape
  >{\centering\arraybackslash}p{1.0cm}     % Route
  >{\raggedright\arraybackslash}p{1.6cm}   % Normalization
  >{\raggedright\arraybackslash}p{1.7cm}   % Augmentation
  >{\centering\arraybackslash}p{1.1cm}     % Overfit Risk
  >{\centering\arraybackslash}p{1.1cm}     % Interpretability
  >{\centering\arraybackslash}p{1.0cm}     % Edge Suit.
  >{\centering\arraybackslash}p{1.0cm}     % Noise Rob.
  >{\raggedright\arraybackslash}p{2.8cm}   % Key Failure Mode
}
\caption{Qualitative comparison of candidate algorithms for three-class $S_{11}$
         classification (\texttt{dry}/\texttt{wet}/\texttt{ice}).
         Assessments reflect expected behaviour given the dataset size
         ($N=1603$) and are \emph{not} derived from experimental results.}
\label{tab:algorithm_comparison} \\

\toprule
\textbf{Algorithm} &
\textbf{Family} &
\textbf{Input Shape} &
\textbf{Route} &
\textbf{Preferred Norm.} &
\textbf{Augment.\ Compat.} &
\textbf{Overfit Risk} &
\textbf{Interp.} &
\textbf{Edge Suit.} &
\textbf{Noise Rob.} &
\textbf{Key Failure Mode} \\
\midrule
\endfirsthead

% ── continuation header ──
\multicolumn{11}{l}{\small\itshape (continued from previous page)} \\[2pt]
\toprule
\textbf{Algorithm} &
\textbf{Family} &
\textbf{Input Shape} &
\textbf{Route} &
\textbf{Preferred Norm.} &
\textbf{Augment.\ Compat.} &
\textbf{Overfit Risk} &
\textbf{Interp.} &
\textbf{Edge Suit.} &
\textbf{Noise Rob.} &
\textbf{Key Failure Mode} \\
\midrule
\endhead

% ── footer ──
\midrule
\multicolumn{11}{r}{\small\itshape (continued on next page)} \\
\endfoot

% ── last footer ──
\bottomrule
\endlastfoot

% ═══════════════════════════════════════════════════════════════
%  ROW 1 — Logistic Regression
% ═══════════════════════════════════════════════════════════════
Logistic Regression
  & Linear
  & $(D,)$ flat vector
  & A
  & StandardScaler (zero-mean, unit-var)
  & Limited; noise injection only
  & Low
  & High
  & High
  & Low
  & Under-fits if class boundaries are non-linear in spectral space \\

% ═══════════════════════════════════════════════════════════════
%  ROW 2 — Linear SVM
% ═══════════════════════════════════════════════════════════════
Linear SVM
  & Linear
  & $(D,)$ flat vector
  & A
  & StandardScaler
  & Limited; noise injection only
  & Low
  & High
  & High
  & Low--Med
  & Same linear limitation; sensitive to outlier support vectors \\

% ═══════════════════════════════════════════════════════════════
%  ROW 3 — RBF SVM
% ═══════════════════════════════════════════════════════════════
RBF SVM
  & Kernel
  & $(D,)$ flat vector
  & A
  & StandardScaler (critical)
  & Limited; noise injection, SMOTE
  & Medium
  & Low
  & Medium
  & Medium
  & $\gamma$--$C$ grid search costly; kernel matrix scales $O(N^2)$ \\

% ═══════════════════════════════════════════════════════════════
%  ROW 4 — Random Forest
% ═══════════════════════════════════════════════════════════════
Random Forest
  & Ensemble (bagging)
  & $(D,)$ flat vector
  & A
  & Not required (tree-based)
  & SMOTE, noise injection
  & Low--Med
  & Medium
  & High
  & High
  & Can memorise noisy features if trees are too deep; feature importance may
    mislead with correlated spectral bins \\

% ═══════════════════════════════════════════════════════════════
%  ROW 5 — Gradient Boosted Trees (XGBoost)
% ═══════════════════════════════════════════════════════════════
Gradient Boosted Trees (XGBoost)
  & Ensemble (boosting)
  & $(D,)$ flat vector
  & A
  & Not required (tree-based)
  & SMOTE, noise injection
  & Medium
  & Medium
  & High
  & High
  & Prone to over-fitting with many boosting rounds on small $N$;
    requires careful early stopping \\

% ═══════════════════════════════════════════════════════════════
%  ROW 6 — Multilayer Perceptron
% ═══════════════════════════════════════════════════════════════
Multilayer Perceptron
  & Neural network (feedforward)
  & $(D,)$ flat vector
  & A
  & StandardScaler or MinMax
  & Noise injection, mixup, dropout
  & High
  & Low
  & Medium
  & Medium
  & Over-parameterised for small $N$; sensitive to learning rate and
    architecture choices \\

% ═══════════════════════════════════════════════════════════════
%  ROW 7 — 1D CNN
% ═══════════════════════════════════════════════════════════════
1D CNN
  & Neural network (convolutional)
  & $(L, C)$ sequence
  & A
  & Per-channel StandardScaler or BatchNorm
  & Noise injection, frequency masking, mixup
  & High
  & Low
  & Medium
  & Med--High
  & Needs sufficient depth for receptive field; over-fits quickly on 1603
    samples without strong regularisation \\

% ═══════════════════════════════════════════════════════════════
%  ROW 8 — ResNet-1D
% ═══════════════════════════════════════════════════════════════
ResNet-1D
  & Neural network (residual)
  & $(L, C)$ sequence
  & A
  & BatchNorm (built-in)
  & Noise injection, frequency masking, mixup
  & High
  & Low
  & Low--Med
  & Med--High
  & Highest parameter count; extremely prone to over-fitting at $N=1603$;
    large model footprint challenges edge deployment \\

% ═══════════════════════════════════════════════════════════════
%  ROW 9 — KAN
% ═══════════════════════════════════════════════════════════════
Kolmogorov--Arnold Network (KAN)
  & Neural network (spline-based)
  & $(D,)$ flat vector
  & A
  & StandardScaler
  & Noise injection, mixup
  & Med--High
  & Medium
  & Medium
  & Medium
  & Immature tooling; spline knot count acts as hidden capacity control;
    limited edge-runtime support \\

% ═══════════════════════════════════════════════════════════════
%  ROW 10 — K-Means (unsupervised)
% ═══════════════════════════════════════════════════════════════
K-Means\textsuperscript{\dag}
  & Unsupervised (centroid)
  & $(D,)$ flat vector
  & A
  & MinMaxScaler or StandardScaler
  & N/A (unsupervised)
  & Low
  & Medium
  & High
  & Low
  & Assumes spherical clusters; label assignment is post-hoc; cannot capture
    non-convex class boundaries \\

\end{longtable}

\vspace{0.3em}
\noindent\textsuperscript{\dag}\,K-Means is included as an \emph{unsupervised probe}
to assess natural cluster separability in the feature space, not as a
direct classification algorithm.  Cluster-to-class mapping is evaluated
via majority voting and the adjusted Rand index.

\end{footnotesize}
 from a section file.
% Designed for landscape orientation; wrap the \input{} call inside
% a \begin{landscape}...\end{landscape} environment (from pdflscape) or
% place it on a standalone landscape page.

\begin{footnotesize}
\setlength{\LTcapwidth}{\linewidth}
\setlength{\tabcolsep}{3.5pt}
\renewcommand{\arraystretch}{1.25}

\begin{longtable}{%
  >{\raggedright\arraybackslash}p{2.1cm}   % Algorithm
  >{\raggedright\arraybackslash}p{1.5cm}   % Family
  >{\raggedright\arraybackslash}p{1.8cm}   % Input Shape
  >{\centering\arraybackslash}p{1.0cm}     % Route
  >{\raggedright\arraybackslash}p{1.6cm}   % Normalization
  >{\raggedright\arraybackslash}p{1.7cm}   % Augmentation
  >{\centering\arraybackslash}p{1.1cm}     % Overfit Risk
  >{\centering\arraybackslash}p{1.1cm}     % Interpretability
  >{\centering\arraybackslash}p{1.0cm}     % Edge Suit.
  >{\centering\arraybackslash}p{1.0cm}     % Noise Rob.
  >{\raggedright\arraybackslash}p{2.8cm}   % Key Failure Mode
}
\caption{Qualitative comparison of candidate algorithms for three-class $S_{11}$
         classification (\texttt{dry}/\texttt{wet}/\texttt{ice}).
         Assessments reflect expected behaviour given the dataset size
         ($N=1603$) and are \emph{not} derived from experimental results.}
\label{tab:algorithm_comparison} \\

\toprule
\textbf{Algorithm} &
\textbf{Family} &
\textbf{Input Shape} &
\textbf{Route} &
\textbf{Preferred Norm.} &
\textbf{Augment.\ Compat.} &
\textbf{Overfit Risk} &
\textbf{Interp.} &
\textbf{Edge Suit.} &
\textbf{Noise Rob.} &
\textbf{Key Failure Mode} \\
\midrule
\endfirsthead

% ── continuation header ──
\multicolumn{11}{l}{\small\itshape (continued from previous page)} \\[2pt]
\toprule
\textbf{Algorithm} &
\textbf{Family} &
\textbf{Input Shape} &
\textbf{Route} &
\textbf{Preferred Norm.} &
\textbf{Augment.\ Compat.} &
\textbf{Overfit Risk} &
\textbf{Interp.} &
\textbf{Edge Suit.} &
\textbf{Noise Rob.} &
\textbf{Key Failure Mode} \\
\midrule
\endhead

% ── footer ──
\midrule
\multicolumn{11}{r}{\small\itshape (continued on next page)} \\
\endfoot

% ── last footer ──
\bottomrule
\endlastfoot

% ═══════════════════════════════════════════════════════════════
%  ROW 1 — Logistic Regression
% ═══════════════════════════════════════════════════════════════
Logistic Regression
  & Linear
  & $(D,)$ flat vector
  & A
  & StandardScaler (zero-mean, unit-var)
  & Limited; noise injection only
  & Low
  & High
  & High
  & Low
  & Under-fits if class boundaries are non-linear in spectral space \\

% ═══════════════════════════════════════════════════════════════
%  ROW 2 — Linear SVM
% ═══════════════════════════════════════════════════════════════
Linear SVM
  & Linear
  & $(D,)$ flat vector
  & A
  & StandardScaler
  & Limited; noise injection only
  & Low
  & High
  & High
  & Low--Med
  & Same linear limitation; sensitive to outlier support vectors \\

% ═══════════════════════════════════════════════════════════════
%  ROW 3 — RBF SVM
% ═══════════════════════════════════════════════════════════════
RBF SVM
  & Kernel
  & $(D,)$ flat vector
  & A
  & StandardScaler (critical)
  & Limited; noise injection, SMOTE
  & Medium
  & Low
  & Medium
  & Medium
  & $\gamma$--$C$ grid search costly; kernel matrix scales $O(N^2)$ \\

% ═══════════════════════════════════════════════════════════════
%  ROW 4 — Random Forest
% ═══════════════════════════════════════════════════════════════
Random Forest
  & Ensemble (bagging)
  & $(D,)$ flat vector
  & A
  & Not required (tree-based)
  & SMOTE, noise injection
  & Low--Med
  & Medium
  & High
  & High
  & Can memorise noisy features if trees are too deep; feature importance may
    mislead with correlated spectral bins \\

% ═══════════════════════════════════════════════════════════════
%  ROW 5 — Gradient Boosted Trees (XGBoost)
% ═══════════════════════════════════════════════════════════════
Gradient Boosted Trees (XGBoost)
  & Ensemble (boosting)
  & $(D,)$ flat vector
  & A
  & Not required (tree-based)
  & SMOTE, noise injection
  & Medium
  & Medium
  & High
  & High
  & Prone to over-fitting with many boosting rounds on small $N$;
    requires careful early stopping \\

% ═══════════════════════════════════════════════════════════════
%  ROW 6 — Multilayer Perceptron
% ═══════════════════════════════════════════════════════════════
Multilayer Perceptron
  & Neural network (feedforward)
  & $(D,)$ flat vector
  & A
  & StandardScaler or MinMax
  & Noise injection, mixup, dropout
  & High
  & Low
  & Medium
  & Medium
  & Over-parameterised for small $N$; sensitive to learning rate and
    architecture choices \\

% ═══════════════════════════════════════════════════════════════
%  ROW 7 — 1D CNN
% ═══════════════════════════════════════════════════════════════
1D CNN
  & Neural network (convolutional)
  & $(L, C)$ sequence
  & A
  & Per-channel StandardScaler or BatchNorm
  & Noise injection, frequency masking, mixup
  & High
  & Low
  & Medium
  & Med--High
  & Needs sufficient depth for receptive field; over-fits quickly on 1603
    samples without strong regularisation \\

% ═══════════════════════════════════════════════════════════════
%  ROW 8 — ResNet-1D
% ═══════════════════════════════════════════════════════════════
ResNet-1D
  & Neural network (residual)
  & $(L, C)$ sequence
  & A
  & BatchNorm (built-in)
  & Noise injection, frequency masking, mixup
  & High
  & Low
  & Low--Med
  & Med--High
  & Highest parameter count; extremely prone to over-fitting at $N=1603$;
    large model footprint challenges edge deployment \\

% ═══════════════════════════════════════════════════════════════
%  ROW 9 — KAN
% ═══════════════════════════════════════════════════════════════
Kolmogorov--Arnold Network (KAN)
  & Neural network (spline-based)
  & $(D,)$ flat vector
  & A
  & StandardScaler
  & Noise injection, mixup
  & Med--High
  & Medium
  & Medium
  & Medium
  & Immature tooling; spline knot count acts as hidden capacity control;
    limited edge-runtime support \\

% ═══════════════════════════════════════════════════════════════
%  ROW 10 — K-Means (unsupervised)
% ═══════════════════════════════════════════════════════════════
K-Means\textsuperscript{\dag}
  & Unsupervised (centroid)
  & $(D,)$ flat vector
  & A
  & MinMaxScaler or StandardScaler
  & N/A (unsupervised)
  & Low
  & Medium
  & High
  & Low
  & Assumes spherical clusters; label assignment is post-hoc; cannot capture
    non-convex class boundaries \\

\end{longtable}

\vspace{0.3em}
\noindent\textsuperscript{\dag}\,K-Means is included as an \emph{unsupervised probe}
to assess natural cluster separability in the feature space, not as a
direct classification algorithm.  Cluster-to-class mapping is evaluated
via majority voting and the adjusted Rand index.

\end{footnotesize}
 from a section file.
% Designed for landscape orientation; wrap the \input{} call inside
% a \begin{landscape}...\end{landscape} environment (from pdflscape) or
% place it on a standalone landscape page.

\begin{footnotesize}
\setlength{\LTcapwidth}{\linewidth}
\setlength{\tabcolsep}{3.5pt}
\renewcommand{\arraystretch}{1.25}

\begin{longtable}{%
  >{\raggedright\arraybackslash}p{2.1cm}   % Algorithm
  >{\raggedright\arraybackslash}p{1.5cm}   % Family
  >{\raggedright\arraybackslash}p{1.8cm}   % Input Shape
  >{\centering\arraybackslash}p{1.0cm}     % Route
  >{\raggedright\arraybackslash}p{1.6cm}   % Normalization
  >{\raggedright\arraybackslash}p{1.7cm}   % Augmentation
  >{\centering\arraybackslash}p{1.1cm}     % Overfit Risk
  >{\centering\arraybackslash}p{1.1cm}     % Interpretability
  >{\centering\arraybackslash}p{1.0cm}     % Edge Suit.
  >{\centering\arraybackslash}p{1.0cm}     % Noise Rob.
  >{\raggedright\arraybackslash}p{2.8cm}   % Key Failure Mode
}
\caption{Qualitative comparison of candidate algorithms for three-class $S_{11}$
         classification (\texttt{dry}/\texttt{wet}/\texttt{ice}).
         Assessments reflect expected behaviour given the dataset size
         ($N=1603$) and are \emph{not} derived from experimental results.}
\label{tab:algorithm_comparison} \\

\toprule
\textbf{Algorithm} &
\textbf{Family} &
\textbf{Input Shape} &
\textbf{Route} &
\textbf{Preferred Norm.} &
\textbf{Augment.\ Compat.} &
\textbf{Overfit Risk} &
\textbf{Interp.} &
\textbf{Edge Suit.} &
\textbf{Noise Rob.} &
\textbf{Key Failure Mode} \\
\midrule
\endfirsthead

% ── continuation header ──
\multicolumn{11}{l}{\small\itshape (continued from previous page)} \\[2pt]
\toprule
\textbf{Algorithm} &
\textbf{Family} &
\textbf{Input Shape} &
\textbf{Route} &
\textbf{Preferred Norm.} &
\textbf{Augment.\ Compat.} &
\textbf{Overfit Risk} &
\textbf{Interp.} &
\textbf{Edge Suit.} &
\textbf{Noise Rob.} &
\textbf{Key Failure Mode} \\
\midrule
\endhead

% ── footer ──
\midrule
\multicolumn{11}{r}{\small\itshape (continued on next page)} \\
\endfoot

% ── last footer ──
\bottomrule
\endlastfoot

% ═══════════════════════════════════════════════════════════════
%  ROW 1 — Logistic Regression
% ═══════════════════════════════════════════════════════════════
Logistic Regression
  & Linear
  & $(D,)$ flat vector
  & A
  & StandardScaler (zero-mean, unit-var)
  & Limited; noise injection only
  & Low
  & High
  & High
  & Low
  & Under-fits if class boundaries are non-linear in spectral space \\

% ═══════════════════════════════════════════════════════════════
%  ROW 2 — Linear SVM
% ═══════════════════════════════════════════════════════════════
Linear SVM
  & Linear
  & $(D,)$ flat vector
  & A
  & StandardScaler
  & Limited; noise injection only
  & Low
  & High
  & High
  & Low--Med
  & Same linear limitation; sensitive to outlier support vectors \\

% ═══════════════════════════════════════════════════════════════
%  ROW 3 — RBF SVM
% ═══════════════════════════════════════════════════════════════
RBF SVM
  & Kernel
  & $(D,)$ flat vector
  & A
  & StandardScaler (critical)
  & Limited; noise injection, SMOTE
  & Medium
  & Low
  & Medium
  & Medium
  & $\gamma$--$C$ grid search costly; kernel matrix scales $O(N^2)$ \\

% ═══════════════════════════════════════════════════════════════
%  ROW 4 — Random Forest
% ═══════════════════════════════════════════════════════════════
Random Forest
  & Ensemble (bagging)
  & $(D,)$ flat vector
  & A
  & Not required (tree-based)
  & SMOTE, noise injection
  & Low--Med
  & Medium
  & High
  & High
  & Can memorise noisy features if trees are too deep; feature importance may
    mislead with correlated spectral bins \\

% ═══════════════════════════════════════════════════════════════
%  ROW 5 — Gradient Boosted Trees (XGBoost)
% ═══════════════════════════════════════════════════════════════
Gradient Boosted Trees (XGBoost)
  & Ensemble (boosting)
  & $(D,)$ flat vector
  & A
  & Not required (tree-based)
  & SMOTE, noise injection
  & Medium
  & Medium
  & High
  & High
  & Prone to over-fitting with many boosting rounds on small $N$;
    requires careful early stopping \\

% ═══════════════════════════════════════════════════════════════
%  ROW 6 — Multilayer Perceptron
% ═══════════════════════════════════════════════════════════════
Multilayer Perceptron
  & Neural network (feedforward)
  & $(D,)$ flat vector
  & A
  & StandardScaler or MinMax
  & Noise injection, mixup, dropout
  & High
  & Low
  & Medium
  & Medium
  & Over-parameterised for small $N$; sensitive to learning rate and
    architecture choices \\

% ═══════════════════════════════════════════════════════════════
%  ROW 7 — 1D CNN
% ═══════════════════════════════════════════════════════════════
1D CNN
  & Neural network (convolutional)
  & $(L, C)$ sequence
  & A
  & Per-channel StandardScaler or BatchNorm
  & Noise injection, frequency masking, mixup
  & High
  & Low
  & Medium
  & Med--High
  & Needs sufficient depth for receptive field; over-fits quickly on 1603
    samples without strong regularisation \\

% ═══════════════════════════════════════════════════════════════
%  ROW 8 — ResNet-1D
% ═══════════════════════════════════════════════════════════════
ResNet-1D
  & Neural network (residual)
  & $(L, C)$ sequence
  & A
  & BatchNorm (built-in)
  & Noise injection, frequency masking, mixup
  & High
  & Low
  & Low--Med
  & Med--High
  & Highest parameter count; extremely prone to over-fitting at $N=1603$;
    large model footprint challenges edge deployment \\

% ═══════════════════════════════════════════════════════════════
%  ROW 9 — KAN
% ═══════════════════════════════════════════════════════════════
Kolmogorov--Arnold Network (KAN)
  & Neural network (spline-based)
  & $(D,)$ flat vector
  & A
  & StandardScaler
  & Noise injection, mixup
  & Med--High
  & Medium
  & Medium
  & Medium
  & Immature tooling; spline knot count acts as hidden capacity control;
    limited edge-runtime support \\

% ═══════════════════════════════════════════════════════════════
%  ROW 10 — K-Means (unsupervised)
% ═══════════════════════════════════════════════════════════════
K-Means\textsuperscript{\dag}
  & Unsupervised (centroid)
  & $(D,)$ flat vector
  & A
  & MinMaxScaler or StandardScaler
  & N/A (unsupervised)
  & Low
  & Medium
  & High
  & Low
  & Assumes spherical clusters; label assignment is post-hoc; cannot capture
    non-convex class boundaries \\

\end{longtable}

\vspace{0.3em}
\noindent\textsuperscript{\dag}\,K-Means is included as an \emph{unsupervised probe}
to assess natural cluster separability in the feature space, not as a
direct classification algorithm.  Cluster-to-class mapping is evaluated
via majority voting and the adjusted Rand index.

\end{footnotesize}
 from a section file.
% Designed for landscape orientation; wrap the \input{} call inside
% a \begin{landscape}...\end{landscape} environment (from pdflscape) or
% place it on a standalone landscape page.

\begin{footnotesize}
\setlength{\LTcapwidth}{\linewidth}
\setlength{\tabcolsep}{3.5pt}
\renewcommand{\arraystretch}{1.25}

\begin{longtable}{%
  >{\raggedright\arraybackslash}p{2.1cm}   % Algorithm
  >{\raggedright\arraybackslash}p{1.5cm}   % Family
  >{\raggedright\arraybackslash}p{1.8cm}   % Input Shape
  >{\centering\arraybackslash}p{1.0cm}     % Route
  >{\raggedright\arraybackslash}p{1.6cm}   % Normalization
  >{\raggedright\arraybackslash}p{1.7cm}   % Augmentation
  >{\centering\arraybackslash}p{1.1cm}     % Overfit Risk
  >{\centering\arraybackslash}p{1.1cm}     % Interpretability
  >{\centering\arraybackslash}p{1.0cm}     % Edge Suit.
  >{\centering\arraybackslash}p{1.0cm}     % Noise Rob.
  >{\raggedright\arraybackslash}p{2.8cm}   % Key Failure Mode
}
\caption{Qualitative comparison of candidate algorithms for three-class $S_{11}$
         classification (\texttt{dry}/\texttt{wet}/\texttt{ice}).
         Assessments reflect expected behaviour given the dataset size
         ($N=1603$) and are \emph{not} derived from experimental results.}
\label{tab:algorithm_comparison} \\

\toprule
\textbf{Algorithm} &
\textbf{Family} &
\textbf{Input Shape} &
\textbf{Route} &
\textbf{Preferred Norm.} &
\textbf{Augment.\ Compat.} &
\textbf{Overfit Risk} &
\textbf{Interp.} &
\textbf{Edge Suit.} &
\textbf{Noise Rob.} &
\textbf{Key Failure Mode} \\
\midrule
\endfirsthead

% ── continuation header ──
\multicolumn{11}{l}{\small\itshape (continued from previous page)} \\[2pt]
\toprule
\textbf{Algorithm} &
\textbf{Family} &
\textbf{Input Shape} &
\textbf{Route} &
\textbf{Preferred Norm.} &
\textbf{Augment.\ Compat.} &
\textbf{Overfit Risk} &
\textbf{Interp.} &
\textbf{Edge Suit.} &
\textbf{Noise Rob.} &
\textbf{Key Failure Mode} \\
\midrule
\endhead

% ── footer ──
\midrule
\multicolumn{11}{r}{\small\itshape (continued on next page)} \\
\endfoot

% ── last footer ──
\bottomrule
\endlastfoot

% ═══════════════════════════════════════════════════════════════
%  ROW 1 — Logistic Regression
% ═══════════════════════════════════════════════════════════════
Logistic Regression
  & Linear
  & $(D,)$ flat vector
  & A
  & StandardScaler (zero-mean, unit-var)
  & Limited; noise injection only
  & Low
  & High
  & High
  & Low
  & Under-fits if class boundaries are non-linear in spectral space \\

% ═══════════════════════════════════════════════════════════════
%  ROW 2 — Linear SVM
% ═══════════════════════════════════════════════════════════════
Linear SVM
  & Linear
  & $(D,)$ flat vector
  & A
  & StandardScaler
  & Limited; noise injection only
  & Low
  & High
  & High
  & Low--Med
  & Same linear limitation; sensitive to outlier support vectors \\

% ═══════════════════════════════════════════════════════════════
%  ROW 3 — RBF SVM
% ═══════════════════════════════════════════════════════════════
RBF SVM
  & Kernel
  & $(D,)$ flat vector
  & A
  & StandardScaler (critical)
  & Limited; noise injection, SMOTE
  & Medium
  & Low
  & Medium
  & Medium
  & $\gamma$--$C$ grid search costly; kernel matrix scales $O(N^2)$ \\

% ═══════════════════════════════════════════════════════════════
%  ROW 4 — Random Forest
% ═══════════════════════════════════════════════════════════════
Random Forest
  & Ensemble (bagging)
  & $(D,)$ flat vector
  & A
  & Not required (tree-based)
  & SMOTE, noise injection
  & Low--Med
  & Medium
  & High
  & High
  & Can memorise noisy features if trees are too deep; feature importance may
    mislead with correlated spectral bins \\

% ═══════════════════════════════════════════════════════════════
%  ROW 5 — Gradient Boosted Trees (XGBoost)
% ═══════════════════════════════════════════════════════════════
Gradient Boosted Trees (XGBoost)
  & Ensemble (boosting)
  & $(D,)$ flat vector
  & A
  & Not required (tree-based)
  & SMOTE, noise injection
  & Medium
  & Medium
  & High
  & High
  & Prone to over-fitting with many boosting rounds on small $N$;
    requires careful early stopping \\

% ═══════════════════════════════════════════════════════════════
%  ROW 6 — Multilayer Perceptron
% ═══════════════════════════════════════════════════════════════
Multilayer Perceptron
  & Neural network (feedforward)
  & $(D,)$ flat vector
  & A
  & StandardScaler or MinMax
  & Noise injection, mixup, dropout
  & High
  & Low
  & Medium
  & Medium
  & Over-parameterised for small $N$; sensitive to learning rate and
    architecture choices \\

% ═══════════════════════════════════════════════════════════════
%  ROW 7 — 1D CNN
% ═══════════════════════════════════════════════════════════════
1D CNN
  & Neural network (convolutional)
  & $(L, C)$ sequence
  & A
  & Per-channel StandardScaler or BatchNorm
  & Noise injection, frequency masking, mixup
  & High
  & Low
  & Medium
  & Med--High
  & Needs sufficient depth for receptive field; over-fits quickly on 1603
    samples without strong regularisation \\

% ═══════════════════════════════════════════════════════════════
%  ROW 8 — ResNet-1D
% ═══════════════════════════════════════════════════════════════
ResNet-1D
  & Neural network (residual)
  & $(L, C)$ sequence
  & A
  & BatchNorm (built-in)
  & Noise injection, frequency masking, mixup
  & High
  & Low
  & Low--Med
  & Med--High
  & Highest parameter count; extremely prone to over-fitting at $N=1603$;
    large model footprint challenges edge deployment \\

% ═══════════════════════════════════════════════════════════════
%  ROW 9 — KAN
% ═══════════════════════════════════════════════════════════════
Kolmogorov--Arnold Network (KAN)
  & Neural network (spline-based)
  & $(D,)$ flat vector
  & A
  & StandardScaler
  & Noise injection, mixup
  & Med--High
  & Medium
  & Medium
  & Medium
  & Immature tooling; spline knot count acts as hidden capacity control;
    limited edge-runtime support \\

% ═══════════════════════════════════════════════════════════════
%  ROW 10 — K-Means (unsupervised)
% ═══════════════════════════════════════════════════════════════
K-Means\textsuperscript{\dag}
  & Unsupervised (centroid)
  & $(D,)$ flat vector
  & A
  & MinMaxScaler or StandardScaler
  & N/A (unsupervised)
  & Low
  & Medium
  & High
  & Low
  & Assumes spherical clusters; label assignment is post-hoc; cannot capture
    non-convex class boundaries \\

\end{longtable}

\vspace{0.3em}
\noindent\textsuperscript{\dag}\,K-Means is included as an \emph{unsupervised probe}
to assess natural cluster separability in the feature space, not as a
direct classification algorithm.  Cluster-to-class mapping is evaluated
via majority voting and the adjusted Rand index.

\end{footnotesize}
